\subsection{The exact common \texorpdfstring{$A_6$}{A6} and double-cover source and its full boundary fibre}

\begin{statement}{E.54}{Proved}{The full $A_6$ and signed Ferrari refinements coexist with the retained boundary family}
Write $G,C,P,E,T_{A_6}$ for the full $A_6$ source, its
perfect residual, its $N$ image, its $F'$ image, and its
four-strand target in E.48.  Write $G_L,C_L,T_L$ for
the signed source and target of E.38, and put
$\Sigma_L=G_L/C_L$.  All these symbols retain their
specified maps from $B$ or $B_4$.

There is a specified quotient $q_{30}:B\twoheadrightarrow
\mathcal Q_{30}$ with
\[
 \begin{gathered}
 |\mathcal Q_{30}|=1166400000,\qquad
 q_{30}(F)=\mathcal A\times(\mathbf Z/30)^2,\\
 \mathcal A=A_5^3=q_{30}(F'),\qquad
 \mathcal Q_{30}/\mathcal A
       =B/(F'F^{30})=:\mathcal A_{30}.
 \end{gathered}
 \tag{E.54.1}
\]
The three maps to $A_5$ are the source epikernels
$4,12,18$ specified in the proof.  The last two coordinates
are the ordinary $x,y$ exponent sums.  The maps $G\to
\mathcal Q_{30}$ and $G_L\to\mathcal Q_{30}$ are onto,
and the actual joint source is exactly
\[
 J=\operatorname{im}\bigl(B\longrightarrow G\times G_L\bigr)
     =G\times_{\mathcal Q_{30}}G_L.
 \tag{E.54.2}
\]
Its perfect residual and its maximum solvable quotient are
\[
 \begin{gathered}
 C_J=C\times_{\mathcal A}C_L=C_J',\\
 J/C_J\simeq (G/C)\times_{\mathcal A_{30}}\Sigma_L.
 \end{gathered}
 \tag{E.54.3}
\]

Let $R_{\rm nil}(H)=\bigcap_{n\ge1}\gamma_n\widehat H$,
using closed lower-central subgroups, and set
$I=\gamma_5\widehat N\cap R_{\rm nil}(H)$.
The following is an exact image of an actual closed source
subgroup:
\[
 \left.(\rho,\rho_L)\right|_I:
 I\twoheadrightarrow
       P'\times_{\mathcal A}C_L.
 \tag{E.54.4}
\]
In particular $I\subset\widehat F'$, its $D_{23}$ image is
trivial, and it has trivial image in every finite quotient
of $B$ whose $H$ image is nilpotent.

Every one of the $864$ E.24 points, at every compatible
multiplier $\lambda\equiv1\pmod{72}$, admits one commutator
word retaining its full $D_{23}$ coordinate and satisfying
both full Ferrari tests, with
\[
       \rho(f)=\rho(f_\sigma),\qquad
       \rho_L(f)=\rho_L(f_\sigma)
 \tag{E.54.5}
\]
for the same actual Galois element $\sigma$ of character
$\lambda$.  The corrected map is onto the actual common
source in $D_{23}\times G_L\times G$.

For multiplier one, this extension also retains every
prescribed set of boundary values admitted by E.43 and
its E.44 and E.47 extensions.  Thus all $432$ old points,
their retained two-adic and three-adic boundary fibres,
and the identity maps on the full relative pro-$p$ boundary
factors for $p\ge5$ coexist with \textup{(E.54.5)}.
Every resulting finite common source map is onto.
The correction from the chosen old word lies in
$I\cap\ker\rho_L$.

The joint new four-strand target is the actual group
\[
 T_{L,A_6}:=\operatorname{im}
       \bigl(B_4\longrightarrow T_L\times T_{A_6}\bigr).
 \tag{E.54.6}
\]
Its maximum solvable quotient is the specified old
$S_L$ of order $96$, through its projection to $T_L$.
The original and corrected common four-strand images
are both exactly
\[
                  T_{23}\times_{S_L}T_{L,A_6}.
 \tag{E.54.6a}
\]
Thus the joint new map is onto, and every retained old
target value is preserved.  All six evaluations in the
two new factors, their pentagon, braid, full-twist and
periodic-conjugacy identities are those of the same
Galois action.
\end{statement}

\begin{proof}
We first identify the entire common source and its central
obstruction, and then construct corrections in the exact
closed subgroup appearing in \textup{(E.54.4)}.

\paragraph{The actual shared epikernels.}
Use ordinary permutation composition, with the rightmost
factor acting first.  The three specified epimorphisms
$\theta_j:F\to A_5$ are
\[
\begin{array}{c|c|c}
 j&\theta_j(x)&\theta_j(y)\\\hline
 4&(345)&(13245)\\
 12&(12345)&(235)\\
 18&(12345)&(13542).
\end{array}
 \tag{E.54.7}
\]
They are three distinct kernels in the exhaustive $19$-kernel
model of E.29.  Let $b_q$ be the retained representatives of
$B/F$.  On its thirty-point induced block, a source letter
$g$ sends
\[
 (q,i)\longmapsto
 \left(qg,\theta_4(b_qgb_{qg}^{-1})^{-1}(i)\right).
 \tag{E.54.8}
\]
The inverse in this formula converts ordinary permutation
composition to the right-permutation convention of the
certificate.  Its six cosets have the source epikernels
$12,18,4,18,4,12$, respectively.  The pure image is exactly
$\mathcal A=A_5^3$: the distinct simple epikernels give the
full product, and the literal thirty-point order check is
$6\cdot60^3$, with pure order $60^3$.

Every source epikernel is tested against every relation of
the complete source presentations.  The results are
\[
\begin{array}{c|r|l}
 \text{source}&\text{complete relations}&
                    \text{factoring source epikernels}\\\hline
 G&865049&4,12,18\\
 G_A&23945&0,2,4,10,12,16,18.
\end{array}
 \tag{E.54.9}
\]
Each rejected kernel has a saved literal relation word.
An independent NumPy replay checks that word to be the
identity in its full source and nonidentity in the rejected
induced block.  The positive entries in
\textup{(E.54.9)} use the entire presentation, including all
Schreier levels.

For $m=15,30$, adjoin the induced cyclic character
\[
 (q,t)\longmapsto
 \left(qg,t+\operatorname{ab}_x(b_qgb_{qg}^{-1})\pmod m\right).
 \tag{E.54.10}
\]
Its six conjugate characters retain both ordinary exponent
coordinates: the two pure images commute, each has order
$m$, and their generated image has order $m^2$.
The two actual blocks give a faithful model of
$\mathcal Q_m$ on $30+6m$ points, and direct order checks give
\[
 \begin{array}{c|r|r}
 m&|\mathcal Q_m|&|q_m(F)|\\\hline
 15&291600000&48600000\\
 30&1166400000&194400000.
 \end{array}
 \tag{E.54.11}
\]
The pure images are $A_5^3\times(\mathbf Z/m)^2$, through
the specified maps.  The quotient by $A_5^3$ is therefore
the literal ordinary quotient $B/(F'F^m)$; no splitting of
its outer $S_3$ extension is used.

E.48 provides the ordinary exponent quotient modulo $120$
of $G$.  E.29 gives the ordinary exponent quotient modulo
$15$ of $G_A$, and E.43 gives the quotient modulo $30$ of
$G_L$.  Combining these maps with \textup{(E.54.9)} supplies
the indicated maps to $\mathcal Q_{15}$ and
$\mathcal Q_{30}$.  All of them agree on the literal
source generators $s,r$.

The actual joint $G/G_A$ source is compressed to degree
$11010$ with explicit point intertwiners.  Its complete
Schreier chain gives order
\[
                  2^{249}3^{104}5^{78}7^8 11^3 13^2.
 \tag{E.54.12}
\]
Above this actual base, all $1154659$ complete relations
are evaluated in the signed Ferrari model.  Their whole
normal binary image has dimension $42$ in the ambient
$2052$-coordinate module.  Every one of its $42$ basis
words and all $84$ source conjugates are independently
replayed in the full base and signed model.

For comparison with the old $44$-dimensional signed kernel,
the exact E.50 modular matrix is
$I+10\left[\begin{smallmatrix}a&b\\c&a\end{smallmatrix}\right]$
modulo $20$.  Its projective reduction modulo four retains
exactly $b,c$.  The new $42$-space is exactly the kernel of
those two specified rows; $G$ already retains this ordinary
modulo-two exponent quotient.  Thus the comparison includes
the actual signed coordinates.
The resulting joint $G/G_L$ order is
\[
                  2^{291}3^{104}5^{78}7^8 11^3 13^2.
 \tag{E.54.13}
\]
Dividing the product of the two separate source orders by
\textup{(E.54.13)} gives exactly $|\mathcal Q_{30}|$.
The joint source is already a subgroup of the fibre product
over the specified quotient maps; equality of these two
orders proves \textup{(E.54.2)}.  The same comparison with
\textup{(E.54.12)} identifies the maximum common $G/G_A$
quotient with the specified $\mathcal Q_{15}$.

The smaller E.45 source comparison is retained as a control.
Its common quotients with $G_A$ and $G_L$ have orders $1350$
and $5400$, respectively.  The actual full joint source has
index $216000$ in the fibre product obtained using that
smaller source.  Thus the extra factor in the common
comparison is precisely retained by the three actual
epikernel maps above.

\paragraph{The central obstruction and the perfect joint residual.}
Both $C$ and $C_L$ map onto $\mathcal A$ in
\textup{(E.54.1)}.  Indeed their ambient quotients are
solvable, and the image of a perfect group in the solvable
quotient $\mathcal Q_{30}/\mathcal A$ is trivial; the
remaining quotient of $\mathcal A$ must also be solvable.

The map $\operatorname{SL}_2(\mathbf F_5)\to A_5$ used here
is the actual conjugation action on its five Klein-four
subgroups.  Steinberg's universal central-extension theorem
applies to this field; its explicit corollary states that
$\operatorname{SL}_2(\mathbf F_5)\to\operatorname{PSL}_2
(\mathbf F_5)$ is universal
\cite[\S7, Theorem 10 and Corollary 2(b)]{DE_Steinberg1967}.
The product map
\[
 U:=\operatorname{SL}_2(\mathbf F_5)^3
                      \longrightarrow\mathcal A
 \tag{E.54.14}
\]
is universal as well.  In any central extension of
$\mathcal A$, the universal images of its three perfect
factors commute: their cross commutator is central and
factors through their abelianizations.  This proves the
product universal property directly.

Consider the full pullback $\widetilde C=C\times_{\mathcal A}U$.
The ten certified generators of $C$ are expanded into one
literal $F$-word graph; all $22824$ required ancestors are
independently evaluated in the full $9462$-point source.
Their three specified $\operatorname{SL}_2$ values are then
used to evaluate the entire $C$ presentation.  This includes
the $720615$ Schreier relations and all ten original-generator
relations to their normal forms.  The latter relations are
retained because the accepted $C$ chain was constructed in
derived mode.

The abelian relation matrix has ten original-generator
exponent columns and three central columns.  Thirteen actual
relations have the block form
\[
       \begin{pmatrix}I_{10}&A\\0&M\end{pmatrix},
 \qquad
       M=\begin{pmatrix}1&1&0\\0&0&1\\0&1&1\end{pmatrix}
       \quad\text{over }\mathbf F_2.
 \tag{E.54.15}
\]
All thirteen words are independently checked to be the
source identity, and their lifted values are checked by
direct $2\times2$ matrix multiplication.  The determinant
of $M$ is one.  The thirty independent central-sign choices
on the ten generator lifts act by the exact invertible
column changes $(e,z)\mapsto(e,z+eD)$,
$D\in\operatorname{Mat}_{10\times3}(\mathbf F_2)$.
All thirty elementary changes are retained in a separate
certificate.  Consequently \textup{(E.54.15)} covers every
choice of lifts.

Since $C$ is perfect, the abelianization of its central
extension $\widetilde C$ is generated by its exponent-two
center.  The full rank in \textup{(E.54.15)} therefore proves
$\widetilde C'=\widetilde C$ integrally.
To obtain the needed exact kernel comparison, put
$K=\ker(C\to\mathcal A)$ and
$A=C/[C,K]$.  This is a perfect central extension of
$\mathcal A$.  Universality gives $\phi:U\to A$ over
$\mathcal A$.  Its image $V$ satisfies $A=Z(A)V$,
so perfection gives $A=[V,V]\subset V$; thus $\phi$ is onto.
There is an onto homomorphism
\[
 \widetilde C\longrightarrow K/[C,K],
       \qquad(g,u)\longmapsto\bar g\phi(u)^{-1}.
 \tag{E.54.16}
\]
Its value is central, which proves the multiplication rule;
the elements $(g,1)$, $g\in K$, prove surjectivity.
Perfection of $\widetilde C$ makes this abelian image trivial.
Hence $K=[C,K]$.

Now $C_J=C\times_{\mathcal A}C_L$ projects onto $C_L$ with
kernel $K$.  It projects onto $C$ as well, so
$[C_J,K]=[C,K]=K$.  Its derived group projects onto the
perfect group $C_L$ and contains $K$.  Therefore $C_J$ is
perfect.  Quotienting \textup{(E.54.2)} by $C_J$ gives the
displayed fibre product in \textup{(E.54.3)}: lifts of any
compatible pair of solvable coordinates can be adjusted in
$C$, which maps onto $\mathcal A$, to agree in the full
common quotient.  That quotient is solvable, proving the
asserted maximality.

\paragraph{The exact lower-central correction image.}
The signed source retains the actual quotient $Q=B/N$ of
order $576$.  Thus the $N$ image of the actual joint source
is exactly the inverse image of the identity under that
specified quotient.  In $G_L$, the quotient
$\rho_L(N)/C_L$ has order $75$.  Under the ordinary
modulo-thirty exponent map it is exactly
$L/30\mathbf Z^2$, where $L$ is the lattice in E.48.8;
the image has order $75$, so this is an isomorphism with
those specified coordinates.
It follows from \textup{(E.54.2)} that
\[
       \rho_J(N)/C_J\simeq P/C.
 \tag{E.54.17}
\]
The perfect normal subgroup $C_J$ belongs to every
lower-central term of $\rho_J(N)$.  E.48 proves
$\gamma_jP=P'$ for every $j\ge2$.  Applying
\textup{(E.54.17)} therefore gives
\[
       \rho_J(\gamma_j\widehat N)
          =P'\times_{\mathcal A}C_L=:\mathcal R
                       \qquad(j\ge2).
 \tag{E.54.18}
\]

It remains to retain the entire nilpotent boundary family.
The quotient $\mathcal R/C_J$ is $P'/C$, whose retained
ternary and binary layers are $V_6$ and $W$ in E.48.
For the literal $x,y$ action matrices, the new exact checks
give
\[
 \begin{gathered}
    (X^3-1)V_6+(Y^3-1)V_6=V_6,\\
    \sum_{i=1}^{9}(\operatorname{Ad}(e_i)-1)W=W.
 \end{gathered}
 \tag{E.54.19}
\]
Here $x^3,y^3,e_i$ are actual elements of $H$, and $e_i$
are the nine specified $F'$ basis words of E.48.  The
program retains literal $\gamma_5N$ preimages of all six
ternary basis values, obtained by three successive
commutators with the actual $N$ generators.  The binary
inputs are the six literal degree-five words already
checked in E.48.  The two commutator matrices have ranks
six and six; their selected columns and source word graphs
are retained.

Since $C_J$ is perfect and belongs to both $\rho_J(H)$ and
$\mathcal R$, it lies in $[\rho_J(H),\mathcal R]$.
Formula \textup{(E.54.19)} gives the whole quotient by
$C_J$.  Consequently
\[
                  [\rho_J(H),\mathcal R]=\mathcal R.
 \tag{E.54.20}
\]
Starting with $I_0=\gamma_5\widehat N\subset\widehat H$,
put $I_{n+1}=[\widehat H,I_n]$, with closed commutator
subgroups.  Equations \textup{(E.54.18)} and
\textup{(E.54.20)} give $\rho_J(I_n)=\mathcal R$ for every
$n$.  The $I_n$ form a descending sequence and
$I_n\subset\gamma_{n+1}\widehat H\cap\gamma_5\widehat N$.
For each fixed value in the finite group $\mathcal R$,
its nonempty compact preimage fibres in these subgroups
have nonempty intersection.  This proves the surjection
in \textup{(E.54.4)}.  The reverse containment follows
already from \textup{(E.54.18)}.

\paragraph{A single Galois value and all old boundary coordinates.}
For the finite family, begin with the actual E.38 word
$f_0$ and its Galois element $\sigma$, so
$\rho_L(f_0)=\rho_L(f_\sigma)$.
For multiplier one and prescribed boundary values, begin
with the word and the same comparison supplied by E.43,
E.44 and E.47.  Their sign convention is the positive-root
one of E.48.16, and the literal all-$864$ comparison in
E.48.21 identifies the same $\kappa_2\bmod3$ labels.

The exact old-kernel formula E.48.12 and the norm quotient
formula E.48.19 give
$d=\rho(f_0)^{-1}\rho(f_\sigma)\in P'$.
Its $\mathcal A$ coordinate is trivial, because the two
signed values agree and the actual common quotient maps
in \textup{(E.54.2)} agree on every source word.
Thus $(d,1)\in\mathcal R$.
Choose its actual preimage $k\in I$ using
\textup{(E.54.4)} and put $f=f_0k$.
This is a commutator word, retains the signed value, and
has the required $G$ value.  Its old $D_{23}$ value is
unchanged by E.24's $\gamma_5N$ kernel bound.
Every prescribed boundary value is unchanged because
each finite boundary factor has nilpotent $H$ image and
$k\in R_{\rm nil}(H)$.  This proves
\textup{(E.54.5)} simultaneously in all retained factors.

\paragraph{Generation on every actual common source.}
Let $S$ be a finite old solvable source retaining $D_{23}$,
$\Sigma_L$, and the chosen finite nilpotent $H$ boundary
factors.  In the boundary family one takes the cofinal
subsystem retaining $P_{20}$, as in E.43; its specified
quotient to $\Sigma_L$ is then present.  The actual old
source and its corrected image are both
$S\times_{\Sigma_L}G_L$.  For the finite all-$864$ family,
take $S$ to be the actual image in $D_{23}\times\Sigma_L$;
E.38 gives the same assertion.

Let $\mathcal J$ be the actual image of $B$ in
$S\times G_L\times G$.  By \textup{(E.54.4)} its $I$
image is precisely $\mathcal R$, with identity $S$
coordinate.  The corrected image contains this same
subgroup.  To verify that last assertion, the Galois
automorphism $\Phi_\sigma$ preserves $H$, since its
ordinary pure abelianization is multiplication by the
unit $\lambda$.  It preserves $N$ as well: its two literal
images in $Q=B/N$ agree with the retained E.24 cusp-gauge
automorphism, because $G_L$ retains $Q$ and has the same
$f$ value.  Equality on $s,r$ gives this exact quotient
comparison.  Thus $\Phi_\sigma(I)=I$.  The corrected
$G\times G_L$ map is the actual source map precomposed
with this automorphism, while its $S$ map kills $I$.
The $D_{23}$ factor kills $\gamma_5N$, and each boundary
factor, including $\Sigma_L$, kills $R_{\rm nil}(H)$.
Hence its $I$ image is indeed the full $\mathcal R$.

Write $\mathcal J_N$ for the actual $N$ image and $S_N$
for its projection to $S$.  The old corrected map takes
$N$ onto $S_N$: its whole source map is onto and its
actual $Q$ map is the retained automorphism.
After quotienting by $\mathcal R$, the exact projection
maps give an injection
\[
       \mathcal J_N/\mathcal R
                     \longrightarrow S_N\times P/P'.
 \tag{E.54.21}
\]
Indeed the $G_L/C_L$ coordinate is already determined
by $S$, and the remaining kernel is exactly the fibre
product defining $\mathcal R$.
The kernel of the projection in \textup{(E.54.21)} to
$S_N$ is central abelian.  A corrected subgroup onto
$S_N$ is therefore normal, and its possible quotient
is a quotient of that central kernel.

The only possible primes in this quotient are two, three
and five.  At $p=2,3$, the actual abelianization modulo
$p$ of $\mathcal J_N/\mathcal R$ is precisely
$N^{\rm ab}/p$: it is a quotient of that rank-$97$
space and projects onto the full space retained by
$D_{23}$.  E.24 proves that the old corrected map is
invertible on both these actual vector spaces.  Thus
the possible quotient has no nontrivial $C_2$ or $C_3$
quotient.  Its five-primary part is already zero:
the five-primary part of $P/P'$ is the ordinary
$F^{\rm ab}/5=C_5^2$, and this exact coordinate is
retained by $\Sigma_L\subset S$.  The lattice in
E.48.8 has determinant $12$, a unit modulo five,
so this is the full specified coordinate of $N$.
The corrected $N$ image is consequently all of
$\mathcal J_N$.  Finally $\mathcal J/\mathcal J_N=Q$,
and its corrected map is the same onto cusp-gauge
map.  This proves full common-source generation.
The argument uses the central kernel in
\textup{(E.54.21)} and makes no nilpotence assertion
about $S_N$.

\paragraph{The new four-strand comparison.}
The actual free Ferrari image in the joint target is
\[
              A_6^{53}\times
                 \operatorname{SL}_2(\mathbf F_5)^{19}
                         \times C_2^2.
 \tag{E.54.22}
\]
The maps to the first two factors are the specified free
maps in E.48 and E.38.  They have no nontrivial common
quotient: every nontrivial quotient of the first product
has a simple factor $A_6$, whereas the second has only
$A_5$ and $C_2$ as composition factors.  Their joint image
is therefore the full product.  It is perfect, so adjoining
the two actual free modulo-two exponent characters gives
the additional direct factor $C_2^2$ in
\textup{(E.54.22)}.

Kill the perfect normal free subgroup
$K_0=A_6^{53}\times\operatorname{SL}_2(\mathbf F_5)^{19}$
of $T_{L,A_6}$.  The remaining quotient $R$ has a faithful
literal representation on $781=530+251$ points: the first
block is the actual $A_6$ Nielsen action, and the second
is the E.38 signed quotient by its perfect free factor.
In particular the latter block retains the binary cross
terms and the two free characters.  The three original
braid generators, all nine pairwise products, and both
braid relations are checked in these exact coordinates.
The full target order is
\[
              |T_{L,A_6}|=2^{403}3^{163}5^{117}7^8 11^3 13^2.
 \tag{E.54.23}
\]

Five complete Schreier certificates in this $781$-point
quotient have the following indices and relation counts:
\[
 \begin{array}{c|r|r}
 \text{certificate}&[R:U_i]&\text{complete relations}\\\hline
 R&1&421070\\
 U_1&2&491742\\
 U_2&6&543926\\
 U_3&96&377977\\
 U_4&96&399945.
 \end{array}
 \tag{E.54.24}
\]
Here $U_1=R'$.  Each subsequent input is a retained list
of ten actual words in the previous subgroup, and its
certificate gives the complete derived group of those
words.  Every nested original word and every strong
generator is independently replayed.  In particular
$U_4\subset U_3'\subset U_3$ and the last two certified
orders agree, proving $U_3=U_3'$; no generation assertion
is inferred merely from the chosen lists.  A separate
complete check of all $747$ conjugations of the $249$
strong generators of $U_3$ by the three original braid
generators proves $U_3\triangleleft R$.

Restriction to the last $251$ points is the literal old
signed quotient map.  Composing with its retained Cayley
table gives the specified epimorphism $R\to S_L$.
Independent word replay gives the identity in $S_L$ for
every strong generator of $U_3$.  Since $[R:U_3]=96=|S_L|$,
this is its exact kernel.  Its preimage in $T_{L,A_6}$ is
an extension of the perfect group $U_3$ by the perfect
group $K_0$, hence is perfect.  Consequently this
specified $S_L$ is the maximum solvable quotient of the
entire joint target.

Every one of the six evaluations in both full induced
sources equals the evaluation of $f_\sigma$.  The two
standard braid triples therefore jointly equal the
original map $B_4\to T_{L,A_6}$ precomposed with the
actual Galois automorphism of Nakamura--Schneps (1.1).
This proves surjectivity onto the actual image in
\textup{(E.54.6)} and every asserted identity there.
The old target $T_{23}$ is solvable and retains the same
specified $S_L$ quotient by E.38.  Every common quotient
with $T_{L,A_6}$ therefore factors through $S_L$.
Both the original and corrected triples agree on that
exact quotient and are onto both factors.  Goursat's
lemma now gives exactly \textup{(E.54.6a)} for both
images.

The complete replay guide is
\texttt{research/DE\_a6\_gl\_replay.md}; the new program
prefix is \texttt{de\_a6\_gl\_}.  The earlier E.48
certificates remain unchanged.  The complete normal
kernel and central-pullback checks retain all old fibres,
all central signs and all defining relations.
\end{proof}
